\chapter{Parametrising the Surface}\label{ch:dv:parametrising-the-surface}

A market maker's screen shows 400 strikes on 20 expiries for one index, each
with a bid and an ask, several times a second. Behind the screen the desk's
surface has five numbers per expiry. It does not pass through every mid: the
mids are noisy, a quarter of a volatility point here and there, and a curve
that passes through all of them bends between them in ways that create
arbitrage and move the \omterm{def:dv:greeks-and-the-hedging-pnl:greeks}{Greeks} of every option priced off it. It does pass
through every bid--ask band, it cannot produce a negative density, it tells the
desk where the wings go beyond the last quoted strike, and its five numbers can
be watched: when the second one jumps, something happened. This chapter builds
that surface. It introduces the parametrisation used across the industry, its
surface-wide version with guarantees against arbitrage, the theorem that limits
how steep the wings can be, the way to fit to bid and ask rather than to mids,
and the adjustments a surface needs for earnings dates and weekends.

\section{Why parametrise}

Chapter 7 built a surface by interpolating quotes. That works when the quotes are
clean and dense. In practice they are neither: mids carry noise of the order of the
bid--ask spread, far strikes are quoted wide or not at all, and every expiry has its
own set of strikes. A parametric smile does four things an interpolation cannot:
\begin{itemize}
\item it averages the noise across strikes, instead of reproducing it;
\item it can be constrained to be free of \omterm{def:dv:implied-volatility-and-its-surface:static}{static arbitrage};
\item it extrapolates beyond the quoted strikes in a controlled way;
\item it summarises a smile in a few numbers whose day-to-day changes can be monitored
and explained.
\end{itemize}
Choosing the numbers so that the smile matches the quotes is a calibration, in the
sense of One Quant Book 4, chapter 24: a small, usually non-convex, optimisation
problem.

\section{SVI and its surface version}

\begin{definition}[SVI parametrisation]\label{def:dv:parametrising-the-surface:svi}
The \emph{SVI parametrisation}\index{SVI parametrisation} (stochastic volatility
inspired) of a smile writes the \omterm{def:dv:implied-volatility-and-its-surface:surface}{total implied variance} of one expiry as
\[
w(k)=a+b\Bigl(\rho(k-m)+\sqrt{(k-m)^2+\varsigma^2}\Bigr),
\]
with $b\ge0$, $|\rho|<1$, $\varsigma>0$ and $a+b\varsigma\sqrt{1-\rho^2}\ge0$ (so that the
minimum of $w$ is non-negative). The parameters have a shape each: $a$ the level, $b$ the
angle between the wings, $\rho$ the rotation (the skew), $m$ the horizontal position of the
minimum, $\varsigma$ the smoothness of the vertex.
\end{definition}

For $|k|\to\infty$ SVI's total variance is linear in $k$, with slopes $b(1+\rho)$ to the
right and $b(1-\rho)$ to the left. That shape is not an accident: the following theorem says
that no arbitrage-free smile can grow faster, and SVI can reach any slope that is allowed.

\begin{theorem}[Lee's moment formula]\label{thm:dv:parametrising-the-surface:lee}
Let $\beta_R=\limsup_{k\to\infty}w(k)/k$. Then $\beta_R\in[0,2]$, and it is determined by the
number of finite moments of the underlying: with $\tilde p=\sup\{p:\E^{\mathbb Q}S_T^{1+p}<\infty\}$,
\[
\beta_R=2-4\Bigl(\sqrt{\tilde p^2+\tilde p}-\tilde p\Bigr).
\]
The left wing obeys the same law with $\tilde q=\sup\{q:\E^{\mathbb Q}S_T^{-q}<\infty\}$. This is
the \emph{Lee moment formula}\index{Lee moment formula}.
\end{theorem}
\admitted

The bound 2 is reached only if the underlying has no finite moment beyond the first; a
lognormal, with all moments, has slope zero. An SVI fit whose wing slope $b(1\pm\rho)$ exceeds
2 has an arbitrage in its extrapolation, whatever it does inside the quoted range, and a fit
constrained to $b(1+|\rho|)\le2$ cannot produce one there
(\cref{fig:dv:parametrising-the-surface:wings}).

\begin{omfigure}
\begin{tikzpicture}
\begin{axis}[width=10.5cm,height=4.8cm,
  xlabel={log-moneyness $k$},ylabel={total variance $w$},
  tick label style={font=\small},label style={font=\small},
  xmin=-2,xmax=2,ymin=0,ymax=0.6,grid=major,grid style={gray!25},
  legend columns=2,legend style={font=\footnotesize,at={(0.5,-0.3)},anchor=north,draw=none,fill=none,/tikz/every even column/.append style={column sep=8pt}},
  table/col sep=comma]
\addplot[omThm,very thick] table[x=k,y=w]{figdata/derivatives/08-parametrising-the-surface/wings.csv};
\addplot[gray,thick,dashed] table[x=k,y=lee]{figdata/derivatives/08-parametrising-the-surface/wings.csv};
\draw[omDef,thick] (axis cs:-0.40,0) -- (axis cs:-0.40,0.6);
\draw[omDef,thick] (axis cs:0.25,0) -- (axis cs:0.25,0.6);
\node[font=\footnotesize,omDef,anchor=east,align=right] at (axis cs:-0.45,0.48) {quoted range\\between the lines};
\legend{fitted three-month SVI slice,Lee's bound $2|k|$}
\end{axis}
\end{tikzpicture}

\omcaption{The three-month SVI slice of the tutorial extended far beyond its quoted strikes
(between the vertical lines). Its wings are straight, with slopes 0.062 on the left and 0.014 on
the right, far inside Lee's bound $2|k|$ (dashed): the extrapolation cannot create an arbitrage.
Data: the tutorial.}\label{fig:dv:parametrising-the-surface:wings}
\end{omfigure}

A good fit slice by slice does not make a good surface: two slices fitted independently can
cross, which is a \omterm{def:dv:implied-volatility-and-its-surface:static}{calendar arbitrage} (chapter 7). The surface version solves that by tying the
slices together through the at-the-money total variance.

\begin{definition}[Surface SVI]\label{def:dv:parametrising-the-surface:ssvi}
\emph{Surface SVI}\index{surface SVI} (SSVI) writes the whole surface as a function of
\omterm{def:dv:implied-volatility-and-its-surface:surface}{log-moneyness} and the at-the-money total variance $\theta_T$ of each expiry,
\[
w(k,\theta_T)=\frac{\theta_T}2\Bigl(1+\rho\varphi(\theta_T)k+\sqrt{\bigl(\varphi(\theta_T)k+\rho\bigr)^2+1-\rho^2}\Bigr),
\]
with one correlation-like parameter $\rho$ and a curvature function $\varphi$, commonly the
power law $\varphi(\theta)=\eta\theta^{-\gamma}$.
\end{definition}

\begin{theorem}[No-arbitrage conditions for SSVI]\label{thm:dv:parametrising-the-surface:ssvi}
An SSVI surface with $\theta_T$ non-decreasing in $T$ is free of \omterm{def:dv:implied-volatility-and-its-surface:static}{calendar arbitrage} if
$0\le\partial_\theta(\theta\varphi(\theta))\le\frac{1}{\rho^2}\bigl(1+\sqrt{1-\rho^2}\bigr)\varphi(\theta)$;
for the power law this holds whenever $0<\gamma<1$. It is free of \omterm{def:dv:implied-volatility-and-its-surface:static}{butterfly arbitrage} if, for
every $\theta$, $\theta\varphi(\theta)(1+|\rho|)<4$ and $\theta\varphi(\theta)^2(1+|\rho|)\le4$.
\end{theorem}
\admitted

The theorem and its proof are in Gatheral and Jacquier; the first butterfly condition is also
necessary. With the power law and $\gamma=\tfrac12$, the at-the-money skew in volatility is
$\rho\eta/(2\sqrt T)$: the $T^{-1/2}$ decay of equity skews that chapter 12 will revisit.

\begin{example}[One surface, five numbers]\label{ex:dv:parametrising-the-surface:ssvi}
Fitted to the tutorial's five expiries (one month to two years), SSVI returns $\rho=-0.560$,
$\eta=1.095$, $\gamma=0.478$ and satisfies both conditions. It misses the mids by 0.09 volatility
points at one year, 0.76 at three months and 2.0 at one month
(\cref{fig:dv:parametrising-the-surface:ssvi}): the market's short smiles are more skewed than
one $\rho$ can describe. Desks use SSVI as the arbitrage-free backbone and add slice-by-slice
corrections where the quotes demand them.
\end{example}

\begin{omfigure}
\begin{tikzpicture}
\begin{axis}[width=10.5cm,height=5cm,
  xlabel={log-moneyness $k$},ylabel={implied volatility (\%)},
  tick label style={font=\small},label style={font=\small},
  xmin=-0.4,xmax=0.25,ymin=8,ymax=44,grid=major,grid style={gray!25},
  legend columns=3,legend style={font=\footnotesize,at={(0.5,-0.3)},anchor=north,draw=none,fill=none,/tikz/every even column/.append style={column sep=8pt}},
  table/col sep=comma]
\addplot[omProp,only marks,mark=*,mark size=1.3pt] table[x=k,y=m1]{figdata/derivatives/08-parametrising-the-surface/ssvi.csv};
\addplot[omDef,only marks,mark=square*,mark size=1.3pt] table[x=k,y=m3]{figdata/derivatives/08-parametrising-the-surface/ssvi.csv};
\addplot[omThm,only marks,mark=triangle*,mark size=1.6pt] table[x=k,y=y1]{figdata/derivatives/08-parametrising-the-surface/ssvi.csv};
\addplot[omProp,thick] table[x=k,y=m1fit]{figdata/derivatives/08-parametrising-the-surface/ssvi.csv};
\addplot[omDef,thick] table[x=k,y=m3fit]{figdata/derivatives/08-parametrising-the-surface/ssvi.csv};
\addplot[omThm,thick] table[x=k,y=y1fit]{figdata/derivatives/08-parametrising-the-surface/ssvi.csv};
\legend{1-month mids,3-month mids,1-year mids}
\end{axis}
\end{tikzpicture}

\omcaption{One SSVI surface (lines) fitted to five expiries at once, shown at three of them
against the mids (markers). The fit is close at one year and misses the one-month skew, which is
steeper than a single $\rho$ allows. Data: the tutorial.}\label{fig:dv:parametrising-the-surface:ssvi}
\end{omfigure}

\section{Splines and non-parametric fits}

The alternative to a formula is a flexible curve: a cubic spline in $k$ through the quotes, or a
smoothing spline that trades closeness to the quotes against roughness. A spline through the
mids reproduces them exactly, noise included; on the tutorial's three-month slice it strays up to
0.17 volatility points from the smile the quotes were drawn from, twice as far as the SVI fit
(0.08), and its second derivative, which is the density, inherits the noise
(\cref{fig:dv:parametrising-the-surface:fit}). A smoothing spline removes most of that at the cost
of a parameter (the smoothing weight) that has no market meaning; and neither kind of spline says
anything about the wings beyond the last strike, nor guarantees the absence of arbitrage without
explicit constraints. Splines are used where no formula fits (a smile with two humps around a
binary event) and as a check on a parametric fit.

\section{Fitting to bid and ask}

A mid is not a price at which anyone trades. What the market says is that the fair value of each
option lies between its bid and its ask, and any smile inside every band is consistent with it.

\begin{method}[Fitting SVI to a band of quotes]\label{met:dv:parametrising-the-surface:fit}
For one expiry with quotes $w^{\mathrm{bid}}_j\le w^{\mathrm{ask}}_j$ at $k_j$:
\begin{enumerate}
\item Objective: the squared distance by which the model leaves each band, plus a small weight on
the distance to the mid to break ties; weight each strike by the inverse square of its width, so
that tight quotes count more.
\item Inner problem: for fixed $(m,\varsigma)$, with $y=(k-m)/\varsigma$, SVI is linear in
$(a,\ b\rho\varsigma,\ b\varsigma)$; solve by weighted least squares and project on the domain
$b(1+|\rho|)\le2$, $|\rho|<1$, $a\ge0$ (the quasi-explicit method).
\item Outer problem: minimise over the two remaining parameters $(m,\varsigma)$ by a derivative-free
search from a few starting points.
\item Check the result: every fitted point inside its band (or report which are not), the wings within
Lee's bound, the density factor $g$ of chapter 7 non-negative on a fine grid.
\end{enumerate}
\end{method}

On the tutorial's three-month slice, 27 strikes with bands from 0.4 volatility points wide at the
money to 2 points at the far left, the fit lands inside all 27 bands, with a largest distance to a mid
of 0.18 points and parameters $a=0.0019$, $b=0.038$, $\rho=-0.63$, $m=0.059$, $\varsigma=0.066$.

\begin{omfigure}
\begin{tikzpicture}
\begin{axis}[width=10.5cm,height=5cm,
  xlabel={log-moneyness $k$},ylabel={3-month implied volatility (\%)},
  tick label style={font=\small},label style={font=\small},
  xmin=-0.41,xmax=0.26,ymin=10,ymax=37,grid=major,grid style={gray!25},
  legend columns=3,legend style={font=\footnotesize,at={(0.5,-0.3)},anchor=north,draw=none,fill=none,/tikz/every even column/.append style={column sep=8pt}},
  table/col sep=comma]
\addplot[gray,only marks,mark=-,mark size=3pt,error bars/.cd,y dir=both,y explicit]
  table[x=k,y=mid,y error expr=\thisrow{ask}-\thisrow{mid}]{figdata/derivatives/08-parametrising-the-surface/quotes.csv};
\addplot[omThm,very thick] table[x=k,y=svi]{figdata/derivatives/08-parametrising-the-surface/fits.csv};
\addplot[omDef,thin,dashed] table[x=k,y=spline]{figdata/derivatives/08-parametrising-the-surface/fits.csv};
\legend{bid--ask band,SVI fitted to the bands,spline through the mids}
\end{axis}
\end{tikzpicture}

\omcaption{A three-month slice: 27 bid--ask bands, the SVI curve fitted to the bands, and a cubic
spline through the mids. At this scale the two curves coincide; the spline carries the mids' noise
(up to 0.17 points from the true smile) where SVI stays within 0.08. Data: the
tutorial.}\label{fig:dv:parametrising-the-surface:fit}
\end{omfigure}

\section{Events, and business time against calendar time}

A surface assumes that variance accrues smoothly with time. Two facts break that: known events,
which add a lump of variance on one date, and weekends and holidays, on which almost none accrues.

\begin{definition}[Event variance]\label{def:dv:parametrising-the-surface:event}
The \emph{event variance}\index{event variance} of a scheduled announcement (an earnings release, a
central-bank decision, a trial result) is the variance of the log-return caused by the announcement
itself. Between two expiries that bracket the event it is the excess of total variance over what the
diffusive rate of the earlier expiry would accrue:
$w_{\mathrm{event}}=w(T_2)-\bigl(w(T_1)/T_1\bigr)T_2$.
\end{definition}

\begin{example}[An earnings week]\label{ex:dv:parametrising-the-surface:earnings}
A share reports after the close on day 7. At-the-money volatilities are 32.0\% for the expiry in 4
days, 56.1\% in 11 days, 48.2\% in 18 and 44.3\% in 25. The \omterm{def:dv:parametrising-the-surface:event}{event variance} is
$0.561^2\times11/365-0.32^2\times11/365=0.0064$: a standard deviation of 8.0\% for the move on the
announcement, an expected absolute move of $8.0\%\times\sqrt{2/\pi}=6.4\%$. Removing it, the
18-day expiry's volatility is 32.0\%: the whole term structure beyond the event is one diffusive
volatility plus one jump (\cref{fig:dv:parametrising-the-surface:earnings}).
\end{example}

Interpolating total variance linearly in calendar time between the 4-day and 11-day expiries would
spread the jump over the week, and price a 7-day option (expiring before the announcement) at
nearly 50\% instead of 32\%. Studies of individual equity options document the pattern: implied
volatility rises into earnings and drops sharply after the release (Dubinsky and Johannes). Surfaces
therefore carry a calendar of events, each with its own variance, fitted from the term structure
and removed before any interpolation in time.

\begin{omfigure}
\begin{tikzpicture}
\begin{axis}[width=10.5cm,height=4.8cm,
  xlabel={days to expiry},ylabel={at-the-money volatility (\%)},
  tick label style={font=\small},label style={font=\small},
  xmin=0,xmax=27,ymin=25,ymax=62,grid=major,grid style={gray!25},
  legend columns=2,legend style={font=\footnotesize,at={(0.5,-0.3)},anchor=north,draw=none,fill=none,/tikz/every even column/.append style={column sep=8pt}},
  table/col sep=comma]
\addplot[omThm,very thick,mark=*,mark size=2pt] table[x=days,y=vol]{figdata/derivatives/08-parametrising-the-surface/earnings.csv};
\addplot[omDef,thick,dashed] table[x=days,y=base]{figdata/derivatives/08-parametrising-the-surface/earnings.csv};
\draw[gray,dotted,thick] (axis cs:7,25) -- (axis cs:7,62);
\node[font=\footnotesize,anchor=west] at (axis cs:7.3,28) {earnings};
\legend{quoted,ex-event: 32\%}
\end{axis}
\end{tikzpicture}

\omcaption{At-the-money term structure of a share around its earnings release (dotted). Expiries after
the release carry one lump of variance, worth 8.0\% of standard deviation; its weight falls as
$1/T$ with the expiry. Illustrative quotes.}\label{fig:dv:parametrising-the-surface:earnings}
\end{omfigure}

\begin{definition}[Business time]\label{def:dv:parametrising-the-surface:business}
\emph{Business time}\index{business time} is a clock that counts the variance-bearing time between
two dates rather than calendar days: trading days count one, weekends and holidays a fraction (often
zero), and event days one plus their \omterm{def:dv:parametrising-the-surface:event}{event variance}. Volatilities quoted on a 252-trading-day clock
are in business time.
\end{definition}

A one-week option at 20\% on a 252-day clock (five trading days) is worth what a 20.34\% volatility
on the 365-day calendar clock gives over seven days. The difference shows up in \omterm{def:dv:greeks-and-the-hedging-pnl:greeks}{theta}: on the calendar
clock a Friday-to-Monday roll costs two-sevenths of a week's time value, on the business clock almost
nothing, and the loss happens on the trading days instead. A desk that marks in calendar time and
sees the market re-mark every Monday is using the wrong clock.

\section{Tutorial: fitting a live slice}\label{tut:dv:parametrising-the-surface:tut}

\begin{tutorial}
\textbf{Goal.} Fit SVI to bid and ask on one expiry, SSVI to a whole surface, and strip an earnings event
from a term structure. \textbf{End state:} the four figures of the chapter and the numbers of
\cref{ex:dv:parametrising-the-surface:ssvi,ex:dv:parametrising-the-surface:earnings}.
\begin{tutsteps}
\item \textbf{The inner problem}: for fixed $(m,\varsigma)$ SVI is linear in three numbers; solve and
project on the admissible domain.
\omcode{code/firm/svi/firm_svi.py}{58}{69}{The quasi-explicit inner fit of SVI, with Lee's bound on the wings.}\label{lst:dv:parametrising-the-surface:inner}
\item \textbf{The event}: variance in excess of the earlier expiry's rate, and the move it implies.
\omcode{code/firm/svi/firm_svi.py}{130}{139}{Event variance between two expiries, and the implied move.}\label{lst:dv:parametrising-the-surface:event}
\item \textbf{Run} \texttt{dv\_svi.fit\_slice()}, \texttt{ssvi\_fit()}, \texttt{earnings\_problem()}
and \texttt{fig\_svi.py}.
\end{tutsteps}
\textbf{What to change next.} Double every bid--ask width and watch the fitted parameters drift within
the band from day to day; then give the SSVI fit a separate $\rho$ for the first two expiries and
compare the one-month error.
\end{tutorial}

\section{Build: the SVI fitter}\label{bld:dv:parametrising-the-surface:svi}

\begin{build}
\bfield{Purpose} Turns raw quotes into the miniature firm's surface: one SVI slice per expiry, fitted to
the bands, checked, with SSVI as the fallback and extrapolation where quotes are missing.
\bfield{Interface} \texttt{svi(k, a, b, rho, m, s)}; \texttt{svi\_check}; \texttt{fit\_svi(k, w\_mid,
w\_bid, w\_ask, weights) -> (params, rmse)}; \texttt{ssvi(k, theta, rho, eta, gamma)};
\texttt{ssvi\_check}; \texttt{fit\_ssvi(slices)}; \texttt{event\_variance}; \texttt{implied\_move};
\texttt{nelder\_mead}.
\bfield{Rules} Fit in total variance; Lee's bound $b(1+|\rho|)\le2$ imposed in the fit; report the
strikes left outside their band; never pass a slice with a negative density factor downstream.
\bfield{Acceptance tests} \texttt{code/firm/svi/tests/}: an exact SVI slice is recovered to
$10^{-5}$ in total variance; the fitted slice of the tutorial lies inside all its bands on several
seeds; SSVI conditions hold for the power law with $\gamma<1$ and fail when $\theta\varphi$ is too
large; \omterm{def:dv:parametrising-the-surface:event}{event variance} recovers a planted jump.
\bfield{Stretch} Joint fit of SVI slices with calendar-arbitrage penalties between them, and a daily
refit that starts from yesterday's parameters (chapter 24).
\end{build}

\begin{omsources}
\item J.~Gatheral, ``A parsimonious arbitrage-free implied volatility parameterization with
application to the valuation of volatility derivatives'', presentation, Global Derivatives, Madrid
(2004).
\item J.~Gatheral and A.~Jacquier, ``Arbitrage-free SVI volatility surfaces'', \emph{Quantitative
Finance} 14 (2014) 59--71.
\item R.~W.~Lee, ``The moment formula for implied volatility at extreme strikes'', \emph{Mathematical
Finance} 14 (2004) 469--480.
\item C.~Martini and S.~De~Marco, ``Quasi-explicit calibration of Gatheral's SVI model'', Zeliade
white paper (2009, revised 2012).
\item A.~Dubinsky and M.~Johannes, ``Earnings announcements and equity options'', working paper,
Columbia University (2006).
\end{omsources}

\section{Exercises}

\begin{exercise}[$\star$]\label{exo:dv:parametrising-the-surface:1}
For the fitted three-month slice ($a=0.0019$, $b=0.038$, $\rho=-0.63$, $m=0.059$, $\varsigma=0.066$),
give the total variance and the implied volatility at the money ($k=0$).
\end{exercise}

\begin{exercise}[$\star$]\label{exo:dv:parametrising-the-surface:2}
Give the left and right wing slopes of that slice, and the number of finite negative moments that
the left slope implies by Lee's formula.
\end{exercise}

\begin{exercise}[$\star$]\label{exo:dv:parametrising-the-surface:3}
The expiry before an announcement (5 days) is at 30\% and the one after (12 days) at 50\%. Give the \omterm{def:dv:parametrising-the-surface:event}{event
variance}, its standard deviation and the expected absolute move.
\end{exercise}

\begin{exercise}[$\star\star$]\label{exo:dv:parametrising-the-surface:4}
Why is interpolating total variance linearly in calendar time wrong between two expiries that bracket
an earnings release? What does it do to an option expiring the day before the release?
\end{exercise}

\begin{exercise}[$\star\star$]\label{exo:dv:parametrising-the-surface:5}
With the fitted SSVI parameters, check both butterfly conditions at the one-month at-the-money total
variance $\theta=0.00183$.
\end{exercise}

\begin{exercise}[$\star\star$]\label{exo:dv:parametrising-the-surface:6}
A one-week option is quoted at 25\% on a 252-trading-day clock with five trading days left. Give the
equivalent volatility on a 365-day calendar clock over seven days.
\end{exercise}

\begin{exercise}[$\star\star\star$]\label{exo:dv:parametrising-the-surface:7}
\emph{Coding.} Fit SVI to the mids only, with equal weights, and compare the number of strikes inside
their band and the wing slopes with the band fit.
\end{exercise}

\begin{exercise}[$\star\star\star$]\label{exo:dv:parametrising-the-surface:8}
\emph{Find the flaw.} ``The spline passes through every mid, so its fitting error is zero; SVI misses
the mids by up to 0.18 points. The spline is the better surface.''
\end{exercise}

\section{Problem: Earnings Week}

\begin{problem}[{Weekend problem --- the move the options are pricing}]\label{pb:dv:parametrising-the-surface:1}
A share reports earnings after the close on day 7. Its at-the-money implied volatilities, quoted to
0.1 point, are 32.0\% (4 days), 56.1\% (11 days), 48.2\% (18 days) and 44.3\% (25 days). A client asks
how much the market thinks the share will move on the release, and what a straddle bought for the
release costs.

\medskip\noindent\textbf{Part I --- Total variances.}
\begin{enumerate}
\item Give the total variance of the 4-day and the 11-day expiries.
\item What total variance would the 11-day expiry have without the event?
\item Give the \omterm{def:dv:parametrising-the-surface:event}{event variance}.
\item Give its standard deviation and the expected absolute move of a normal jump with that deviation.
\item Why use the 4-day expiry, not the 25-day one, as the diffusive rate?
\end{enumerate}

\medskip\noindent\textbf{Part II --- Consistency.}
\begin{enumerate}[resume]
\item Remove the event from the 18-day expiry: what volatility remains?
\item Do the same for the 25-day expiry.
\item What would a remaining volatility well above 32\% suggest?
\item What does the 11-day volatility become the day after the release, if nothing else changes?
\item Why does the 11-day straddle lose value overnight even if the share moves 6\%?
\end{enumerate}

\medskip\noindent\textbf{Part III --- The surface.}
\begin{enumerate}[resume]
\item How should the surface interpolate between the 4-day and 11-day expiries?
\item Which expiries' smiles should be steeper, and why?
\item How would you treat a weekend between the 4-day and 11-day expiries?
\item A 7-day option expires before the release: what volatility should it carry?
\item What happens to the fitted SVI parameters of the 11-day slice after the release?
\end{enumerate}

\medskip\noindent\textbf{Part IV --- Judgement.}
\begin{enumerate}[resume]
\item Is 6.4\% the market's forecast of the move?
\item How would you compare it with the share's history of earnings moves?
\item What trade expresses the view that the move will be smaller than priced, and what is its risk?
\item State the \emph{named result}: the standard deviation and the expected absolute size of the move
priced by the options.
\item In one sentence: why must events be taken out before interpolating in time?
\end{enumerate}
\end{problem}

\section{Interview questions}

\begin{interviewq}[$\star$ \iqroles{researcher, trader}]\label{iq:dv:parametrising-the-surface:1}
Write SVI and say what each parameter does to the smile.
\end{interviewq}

\begin{interviewq}[$\star$ \iqroles{trader}]\label{iq:dv:parametrising-the-surface:2}
From the at-the-money volatilities of the expiries before and after an earnings release, how do you get
the implied move?
\end{interviewq}

\begin{interviewq}[$\star\star$ \iqroles{researcher}]\label{iq:dv:parametrising-the-surface:3}
Why can \omterm{def:dv:implied-volatility-and-its-surface:surface}{total implied variance} not grow faster than $2|k|$ in the wings?
\end{interviewq}

\begin{interviewq}[$\star\star$ \iqroles{researcher, developer}]\label{iq:dv:parametrising-the-surface:4}
You fit SVI every second on 20 expiries. How do you make the fit fast and stable, and how do you fit to
bid and ask rather than to mids?
\end{interviewq}

\begin{interviewq}[$\star\star$ \iqroles{trader, risk}]\label{iq:dv:parametrising-the-surface:5}
Your desk marks \omterm{def:dv:greeks-and-the-hedging-pnl:greeks}{theta} in calendar time. What happens every Monday morning, and what should change?
\end{interviewq}

\begin{interviewq}[$\star\star\star$ \iqroles{researcher}]\label{iq:dv:parametrising-the-surface:6}
When would you choose SSVI over independent SVI slices, and what do you give up?
\end{interviewq}
